%% kapa: exam=1 round=29 year=2018 period=1 form=A subject=economics subject_ko=경제학원론 seq=24 number=64
%% source: Q-Net 감정평가사 기출문제 (cst00309) artlSeq=5202808 file=extracted/1008/5202808/2018년 제29회 감정평가사 1차시험 기출문제/2018년 제29회 감정평가사 1차시험 기출문제/1교시 A형.pdf
%% conversion: kordoc 4.24.51 -> tools/gen1.py (auto); review=proofread-sonnet; flags=none
\begin{question}{64}
실업에 관한 설명으로 옳지 않은 것은?
\begin{choices}
  \item 균형임금을 초과한 법정 최저임금의 인상은 비자발적 실업을 증가시킨다.
  \item 실업급여 인상과 기간 연장은 자발적 실업 기간을 증가시킨다.
  \item 정부의 확장적 재정정책은 경기적 실업을 감소시킨다.
  \item 인공지능 로봇의 도입은 경기적 실업을 증가시킨다.
  \item 구직자와 구인자의 연결을 촉진하는 정책은 마찰적 실업을 감소시킨다.
\end{choices}
\end{question}
\answer{4}
